\section{Check the published bytes before the next call}
For the worked input, successful decoding returns zero in \texttt{v0} and fifty-two
in \texttt{v1}. At record offsets zero through seven, expect
\texttt{08 07 06 05 04 03 02 01}. These are the counter's little-endian bytes, even
though its incoming CBOR argument was big-endian. At offsets eight through fifteen,
expect \texttt{03 00 00 00 0a 00 00 00}: an identifier length of three and a reading
count of ten, each stored in four bytes.

The identifier occupies offsets sixteen through eighteen and remains
\texttt{00 ff 41}. The first three local samples, at offsets forty-eight through
fifty-three, are \texttt{00 00 17 00 18 00}. Each now has a two-byte slot, including
the values that occupied one byte on the wire. Minus one occupies \texttt{ff ff};
minus 32768, the final sample, occupies \texttt{00 80}. With the supplied sentinel
image, identifier storage at offsets nineteen through forty-seven and sample storage
at offsets sixty-eight through 175 must still contain \texttt{a5}.

Call the bound routine on that record. Its successful result is thirty-seven, which
counts the encoded result rather than the local record or original input. Give the
encoder exactly thirty-seven destination bytes and expect the same result count.
An independent decode of those output bytes must recover the counter, identifier and
ten readings just checked. It will find only three fields: optional key three has
served its validation purpose and is not retained by the application record.

Repeat decoding with a byte appended to the input. This time the decoder returns one
and a zero result count. If the caller record started with sentinel bytes, every one
of its 176 bytes must remain a sentinel. A plausible counter or correct sample prefix
does not compensate for changing that record before the complete input has passed.
